\section{Предпосылки: обучаемые модели и организация вычислений}

\subsection{Аналоговые вычислительные машины 1930-х годов}

Выбор физического способа выполнения операций исследовался ещё
до появления математических моделей нейронных сетей. В 1931 году
В.~Буш описал дифференциальный анализатор Массачусетского
технологического института — механическую машину для решения
обыкновенных дифференциальных уравнений~\cite{bush1931}.
Величины представлялись непрерывными перемещениями и вращениями,
а математические зависимости задавались соединением узлов.

В механическом интеграторе взаимодействие диска и ролика
связывало входное движение с накоплением выходной величины.
Для решения задачи требовалось подготовить соединения,
выбрать масштабы переменных и установить начальные условия.
Вычисление выполнялось работой механизма, а результат
получался в виде кривой. Точность зависела от устройства,
его настройки и процедуры получения результата.

Такой анализатор создавался для инженерных расчётов,
а не для нейросетевых задач. Его значение для истории
вычислительной техники состоит в использовании физических
свойств устройства для выполнения математических операций.
Уже при таком способе вычислений необходимо было согласовывать
сложность задачи, точность механизма и форму представления величин.
Эти вопросы позволяют рассматривать историю аппаратных вычислений
шире истории цифровых процессоров, сохраняя различия между
назначением и устройством машин разных эпох.

\subsection{От логической модели нейрона к обучаемой сети}

В статье 1943 года У.~Маккаллох и У.~Питтс исследовали возможность описывать
нервную активность средствами логического исчисления. Нейрон представлялся
элементом с пороговым условием срабатывания, а сеть — системой связей,
поведение которой можно анализировать формально~\cite{mcculloch1943}.
Историческое значение этой работы состоит в переходе от биологического
объекта к математическому описанию операций. Такое описание позволяло
рассуждать о вычислительных возможностях сети независимо от подробностей
конкретного биологического механизма.

Модель Маккаллоха--Питтса относится к формализации вычислительных
возможностей нейронных сетей. Поведение сети в ней описывается через взаимодействие математически заданных
элементов.
Организация массового выполнения матричных операций и оценка энергии
их исполнения относятся к последующим этапам развития аппаратуры.
Для решения этих задач потребовалось развитие как методов обучения, так и
электронной вычислительной техники.

В работе Ф.~Розенблатта 1958 года перцептрон рассматривался как модель
хранения и организации информации. Наряду с описанием заданной сети рассматривалось изменение её связей в
зависимости от опыта.
Автор связывал обучение с распознаванием образов и обобщением
по совокупности наблюдений~\cite{rosenblatt1958}.
Параметры вычислительного преобразования определяются процедурой
обучения на основании используемых наблюдений.

При оценке обучаемой системы различаются численная точность её параметров и
качество решения целевой задачи. Одинаковое возмущение
параметра может приводить к различным изменениям результатов разных
сетей. В ранних исследованиях перцептрона модель оценивалась по результатам обработки
данных. Методы управления
разрядностью впоследствии использовали зависимость допустимых численных
изменений от свойств модели и критериев её применения.

\subsection{Адаптивные схемы и хранение настраиваемых параметров}

Развитие обучаемых моделей сопровождалось созданием устройств,
в которых параметры изменялись в ходе обучения. В техническом
отчёте Б.~Уидроу 1960 года описан адаптивный линейный элемент
Adaline (Adaptive Linear) и его реализация с электрохимической
памятью~\cite{widrow1960memistor}. Входные сигналы взвешивались
и суммировались, после чего результат сравнивался с порогом.

Adaline и перцептрон представляли разные подходы к обучению
пороговых элементов. В первом случае настройка связывалась
с уменьшением среднеквадратичной ошибки, во втором —
с исправлением ошибочных решений. Это различие позднее
было рассмотрено в обзоре Б.~Уидроу
и М.~Лера 1995 года~\cite{widrow1995}.
Ранние исследования соединяли поиск правила обучения
с созданием схемы, способной автоматически изменять параметры.

Для хранения регулируемого коэффициента использовался
\textit{мемистор} — резистивный элемент, сопротивление которого
изменялось при электрохимическом осаждении металла.
Записанное состояние определяло действие элемента в схеме.
Хранение параметра и его применение при обработке сигнала
были связаны устройством компонента.

Этот пример показывает раннюю постановку задачи хранения
и использования обученных параметров в одной схеме.
Для её реализации требовалось обеспечить точность записи,
сохранение состояния и настройку коэффициентов.
Оценка энергии должна учитывать как обработку сигналов,
так и действия, необходимые для изменения параметров.

\subsection{Тепловыделение и обратимость вычислений}

В работе Р.~Ландауэра 1961 года исследовалась связь логической
необратимости с выделением тепла~\cite{landauer1961}.
Если несколько исходных состояний преобразуются в одно,
по результату нельзя однозначно восстановить исходную информацию.
Примером служит установка бита в ноль независимо от его
предыдущего значения.

Результат Ландауэра связывал минимально необходимое
тепловыделение со стиранием информации при определённых
термодинамических условиях. Тем самым вопрос о затратах
вычисления рассматривался на уровне физических ограничений,
а не только характеристик построенных машин. Для истории
вычислительной техники это означало возможность отдельно
обсуждать неизбежные потери информации и потери,
обусловленные выбранным устройством схемы.

Ч.~Беннетт в 1973 году показал возможность логически обратимого
выполнения вычислений с сохранением информации
о промежуточных шагах~\cite{bennett1973}. Его работа развивала
вопрос о том, насколько затраты зависят от организации
вычислительного процесса. Обратимость стала самостоятельным
предметом исследования, при этом возможность обратимого
логического описания не означала отсутствия потерь
в реальном устройстве.

Эти работы позволили различать ограничения, связанные
с потерей информации, и энергозатраты конкретной реализации.
Для исследования нейросетевой аппаратуры такое различие
задаёт два разных вопроса: какие физические ограничения
действуют при обработке информации и сколько энергии
требуется выбранной системе для выполнения заданной задачи.

\subsection{Комплементарная логика и снижение мощности}

Практические способы снижения мощности развивались
в схемотехнике полупроводниковых устройств. В докладе
Ф.~Уонласса и Ч.-Т.~Са 1963 года была представлена логика
на взаимодополняющих полевых транзисторах с каналами
разного типа~\cite{wanlass1963}. Для таких схем используется
название «комплементарная структура металл--оксид--полупроводник»
(Complementary Metal--Oxide--Semiconductor, CMOS).

В рассмотренном инверторе проводящие состояния транзисторов
чередовались: при открытии одного другой закрывался.
Это позволяло ограничить ток в установившемся состоянии.
Авторы отдельно рассматривали потребление при переключениях
и в режиме ожидания. Такое разделение связывало мощность
не только с количеством элементов, но и с их работой во времени.

Для последующего развития вычислительных устройств
это различие имеет прямое значение. Уменьшение числа
выполняемых операций меняет активность схемы; изменение
размеров и соединений элементов влияет на затраты переключения.
Однако характеристики одной логической ячейки не определяют
энергию процессора вместе с памятью и управляющими цепями.
Ранние работы о логике с малой мощностью потребления позволяют проследить
связь между выбором технологии и организацией вычислений
ещё до появления специализированных нейросетевых кристаллов.

\subsection{Обучение многослойных сетей и рост вычислительных задач}

В статье Д.~Румельхарта, Дж.~Хинтона и Р.~Уильямса 1986 года было показано,
как обратное распространение ошибки позволяет настраивать внутренние
представления многослойной сети. Веса изменяются так, чтобы уменьшить
расхождение между фактическим и требуемым выходом~\cite{rumelhart1986}.
Обучение последовательности слоёв позволяет совместно настраивать
внутренние преобразования и выход модели. Это расширяет состав
операций, участвующих в определении параметров сети.

Вычислительная структура обучения многослойной сети включает
повторное выполнение преобразований векторов, определение производных
и изменение параметров. Применение обученной модели использует
полученные параметры для обработки новых входов. Из-за различий в составе операций и хранении промежуточных данных обучение и
инференс предъявляют разные требования к вычислительной аппаратуре. Публикация
1986 года описывает один из этапов развития метода и его применения для
настройки внутренних представлений сети.

В работе Я.~Лекуна, Л.~Ботту, Й.~Бенжио и П.~Хаффнера 1998 года
градиентное обучение было рассмотрено в составе систем распознавания
документов. Свёрточная сеть использовала структуру изображения,
локальные связи и повторное применение параметров~\cite{lecun1998}.
С точки зрения вычислений это означало регулярное повторение однотипных
действий над разными участками данных. Регулярность этих операций
позволяла распределять их выполнение между параллельно работающими узлами.

Прикладная система распознавания документов также включала выделение
полей, сегментацию и обработку результатов. Её производительность
определялась совокупностью этих этапов. Оценка отдельной арифметической
операции характеризует только часть полного исполнения. Для определения
системного эффекта аппаратного ускорения требуется учитывать долю
ускоряемых действий и затраты остальных этапов обработки документа.

\subsection{Локальная память и пределы параллельного ускорения}

Устройство памяти и возможности параллельного исполнения
исследовались в 1960-х годах для вычислительных систем
общего назначения. М.~Уилкс в статье 1965 года рассматривал
небольшую быструю память, работающую совместно с более
вместительной медленной памятью~\cite{wilkes1965}.
Сначала проверялось наличие требуемого слова в быстрой памяти.
Если оно уже хранилось там, обращение к основной памяти
можно было исключить. Требовалось также согласовывать
копии данных при их изменении.

В такой организации ускорение зависело от повторных
обращений к уже загруженным словам. Поэтому результат
определялся не только временем доступа к каждому уровню,
но и последовательностью запросов программы. Статья была
посвящена организации хранения, а не измерению энергии.
Тем не менее она описывала способ исключать часть
обращений к более удалённой памяти — действие, стоимость
которого впоследствии стало необходимо учитывать
при энергетической оценке нейросетевых систем.

Дж.~Амдал в работе 1967 года исследовал ограничения
ускорения при использовании параллельных вычислителей~\cite{amdahl1967}.
Наряду с параллельно исполняемой работой он учитывал
последовательные действия и управление данными.
Даже ускорение большой части задачи не устраняет
время, затрачиваемое на остальные операции.

Для нейросетевых вычислений этот принцип требует
рассматривать полный процесс исполнения: подготовку входов,
вычисление слоёв, передачу данных и получение ответа.
Он ограничивает выводы, основанные только на скорости
отдельного матричного блока. Оценка ускорения при этом
остаётся отличной от оценки энергии: последняя дополнительно
требует сведений о мощности и длительности работы системы.

\subsection{Параллельные вычисления и сложность обмена данными}

Развитие параллельных вычислений сопровождалось исследованием связи
между количеством арифметических операций и организацией обмена.
Рост числа исполнительных элементов увеличивает требования
к доставке операндов и сохранению промежуточных результатов.
Поэтому структура памяти, локальное повторное использование данных
и разбиение задачи влияют на производительность и энергозатраты аппаратуры.
Эти зависимости сохраняют значение для различных вычислительных
архитектур и классов приложений.

Одним из решений стала систолическая организация, описанная
Х.~Т.~Кунгом в 1982 году~\cite{kung1982}. Регулярно взаимодействующие
элементы выполняют операции над данными, последовательно
передаваемыми между ними. Повторное использование операндов
сокращает требования к внешнему обмену. Систолические структуры
представляют частный способ согласования параллелизма и локальности;
общая задача организации обмена распространяется и на другие
способы исполнения алгоритмов.

В работе Цзя-Вэй Хуна и Х.~Т.~Кунга 1981 года сложность обмена
исследовалась посредством игры с красными и синими камнями~\cite{hong1981}.
Вычисление задаётся графом зависимостей, а состояния значений
соответствуют их размещению в быстрой или медленной памяти.
Для ряда алгоритмов устанавливаются нижние оценки обмена
при ограниченной ёмкости быстрой памяти. Такая постановка
позволила исследовать необходимые перемещения независимо
от конкретного устройства исполнительного узла.

Теоретические оценки обмена определяют ограничения принятой
модели вычислений. Реальная система дополнительно характеризуется
ёмкостью буферов, задержками и пропускной способностью связей.
Стоимость исполнения зависит от распределения данных между
уровнями памяти и частоты повторных обращений. В последующих
исследованиях энергоэффективности эти характеристики учитывались
наряду с количеством арифметических действий.

\subsection{Аналоговая микроэлектроника и нейроморфные системы}

В конце XX века исследовались схемы, использующие
физические свойства электронных компонентов для обработки
сигналов по образцу нервных систем. В статье К.~Мида
1990 года этот подход был изложен применительно
к нейроморфной электронике~\cite{mead1990}.
Автор связывал затраты не только с действием транзистора,
но и с передачей сигналов и организацией всей схемы.

Одним из рассмотренных примеров была кремниевая модель
ранней обработки зрительных сигналов. Локальные связи
между чувствительными элементами и резистивной сетью
использовались для сопоставления освещённости точки
с окружающими участками. Часть обработки выполнялась
непосредственно в аналоговой схеме, рядом с местом
получения сигнала. При таком устройстве вычисления
и передача данных зависели от геометрии соединений.

Мид также рассматривал разброс характеристик транзисторов
и адаптацию для компенсации различий между элементами.
Это связывало выбор физического способа обработки
с требованиями к точности и настройке. Представление сигналов
и способ выполнения операций в аналоговой схеме выбирались
совместно с учётом физических свойств компонентов и соединений.

Работы по нейроморфной электронике показывают,
что исследование энергетических затрат нейроподобных
систем уже в 1990 году включало локальность соединений,
специализацию операций и учёт свойств компонентов.
Такой подход дополняет историю цифровой аппаратуры
и позволяет сравнивать различные способы реализации
нейросетевых вычислений без сведения их к одному типу устройства.

\subsection{Конечная точность как самостоятельная научная проблема}

Другая предпосылка энергоэффективных вычислений связана с численным анализом.
Цифровая машина оперирует конечными представлениями чисел, поэтому
реализация арифметического выражения может зависеть от округления,
диапазона представимых значений и порядка действий. В обзоре Д.~Голдберга
1991 года эти вопросы систематизированы применительно к арифметике
с плавающей точкой~\cite{goldberg1991}. В частности, равенство
алгебраических выражений не гарантирует совпадения их машинных результатов.

Для нейросетевой аппаратуры различаются выполнение вычислений по заранее
заданным численным правилам и изменение арифметики модели. Во втором случае требуется установить допустимость возникающих
расхождений для целевой задачи. Численная близость к результату вычисления
с большей точностью и сохранение прикладных характеристик модели
служат разными критериями оценки. Для их проверки нужны разные эксперименты:
сравнение численных ответов и оценка качества решения задачи на выбранных
данных.

К концу XX века были разработаны обучаемые нейронные сети, способы
параллельного выполнения регулярных операций и методы анализа погрешностей
конечной точности. Эти результаты были получены при решении разных задач. Рост
сложности применяемых моделей и ограничения вычислительной техники потребовали
рассматривать их совместно, в том числе для оценки и сокращения энергии
нейросетевых вычислений.
